\section{Related Work}
\textbf{Group robustness.} Group DRO minimises the worst group loss and needs strong regularisation to generalise \cite{sagawa2020groupdro}. Methods that avoid training-time group labels upweight or contrast examples that a first model gets wrong \cite{liu2021jtt,nam2020lff,zhang2022cnc}, and simple class or group balancing is competitive \cite{idrissi2022balancing}. Benchmarks comparing these methods find that model selection and the availability of group labels strongly affect the ranking \cite{yang2023subpopbench}.

\textbf{Frozen backbones and validation labels.} Deep Feature Reweighting (DFR) retrains the last layer on a group-balanced held-out set and recovers most of the worst-group accuracy of full group-robust training \cite{kirichenko2023dfr}, consistent with the observation that spurious correlations mostly affect the final layer \cite{izmailov2022feature}; later work reduces the annotations needed \cite{labonte2023fewer,qiu2023afr}. DPE \cite{to2025dpe} replaces the linear head with fifteen distance-based prototype heads, trained in sequence on group-balanced subsets of the validation set with a penalty that keeps their prototypes apart, and averages them. In all of these methods the validation group labels are training data, not just a selection signal, so their quality matters directly.

\textbf{Label errors in benchmarks.} Test-set label errors are common in widely used benchmarks and can change which model looks best \cite{northcutt2021pervasive}. For Waterbirds, the dataset repository notes that three species are wrongly marked as waterbirds and states that this does not affect the validity of the benchmark \cite{sagawa2020code}. MaskTune \cite{asgari2022masktune} corrected these labels, attributing them to similar species names, together with images showing several birds, and reported much higher ERM worst-group accuracy afterwards. These works do not separate the effect of the validation labels from that of the test labels, and do not study methods trained on the validation set. We give the exact cause in the generator, count the affected images per split, and measure the effect on DPE and a DFR-style head on both sides.

\textbf{Stacking, cross-fitting and calibration.} Stacking fits a second-level model on out-of-sample predictions \cite{wolpert1992stacking,breiman1996stacked}, and cross-fitting produces each prediction from a model that did not see that example \cite{chernozhukov2018dml}. Platt and temperature scaling fit a monotone map of a score \cite{platt1999probabilistic,guo2017calibration}; for a single binary score such a map changes decisions only through the threshold, which can also be tuned directly \cite{elkan2001cost,menon2021logit}. We use these tools to spend DPE's validation set without holding images back, and test whether they do more than move the threshold.
